\chapter{Connections to Wolfram Metamathematics}\label{ch:wolfram}

Wolfram~\cite{wolfram2022physicalization} introduces the \emph{ruliad} --- the
entangled limit of all possible computations --- and proposes that mathematical
theorems correspond to paths in a metamathematical space.

The universality graph $\UG$ is a concrete instantiation of this idea:
\begin{itemize}
\item Vertices $\leftrightarrow$ computational objects in the ruliad
\item Edges $\leftrightarrow$ constructive proofs (certified paths)
\item Paths $\leftrightarrow$ composed proofs (transitivity)
\item Graph metrics $\leftrightarrow$ metamathematical observables
  (overhead as distance, diameter as proof complexity)
\end{itemize}

The \emph{Principle of Computational Equivalence}~\cite{wolfram2002nks} predicts
that $\UG$ is highly connected: most sufficiently complex systems should be
universal, so the universal core subgraph should have small diameter. The
overhead weights on edges give a quantitative measure of ``how different''
two universal systems are --- a notion that the ruliad framework makes
precise.

Wolfram's program of studying the space of all simple
programs~\cite{wolfram2002nks} motivates the enumeration side: by scanning
small $(s,k)$~TMs, small ECAs, and small tag systems, we populate $\VV$ and
discover new edges. The LLM pipeline automates this exploration.
